\documentclass[12pt,a4paper,openany,oneside]{book}

% --- Paquetes ---
\usepackage[utf8]{inputenc}
\usepackage[spanish, es-noshorthands]{babel}
\usepackage{amsmath, amsfonts, amssymb}
\usepackage{graphicx}
\usepackage{geometry}
\usepackage{hyperref}
\usepackage{booktabs}
\usepackage{fancyhdr}
\usepackage{listings}
\usepackage{xcolor}
\usepackage{tikz}
\usepackage{pgfplots}
\usetikzlibrary{arrows.meta, positioning, shapes.geometric, calc, 3d}
\pgfplotsset{compat=1.18}

\geometry{margin=2.5cm}

% --- Colores y estilos para código Python ---
\definecolor{codegreen}{rgb}{0,0.6,0}
\definecolor{codegray}{rgb}{0.5,0.5,0.5}
\definecolor{codepurple}{rgb}{0.58,0,0.82}
\definecolor{backcolour}{rgb}{0.95,0.95,0.92}

\lstset{
    language=Python,
    backgroundcolor=\color{backcolour},   
    commentstyle=\color{codegreen},
    keywordstyle=\color{blue},
    numberstyle=\tiny\color{codegray},
    stringstyle=\color{codepurple},
    basicstyle=\ttfamily\small,
    breaklines=true,
    frame=single,
    showstringspaces=false
}

% --- Estilo de cabecera y pie ---
\pagestyle{fancy}
\fancyhf{}
\renewcommand{\headrulewidth}{0.4pt}
\renewcommand{\footrulewidth}{0.4pt}
\fancyhead[L]{\small Carlos César Sánchez Coronel}
\fancyhead[R]{\small Sesión 2: Tensores - La Estrella del Data Science}
\fancyfoot[L]{\small Carlos César Sánchez Coronel}
\fancyfoot[C]{\thepage}

% --- Portada ---
\title{
    \textbf{Sesión 2} \\ 
    \Huge \textbf{TENSORES} \\
    \Large Estructuras de datos multidimensionales para IA \\

}
\author{
    \small Por \\ \vspace{0.2cm}
    \Large \textsc{Carlos César Sánchez Coronel}
}
\date{2026}

\begin{document}

\maketitle
\tableofcontents
\addcontentsline{toc}{chapter}{Índice general}

% ------------------------------------------------------------------
% CAPÍTULO 2: SESIÓN 2 - TENSORES
% ------------------------------------------------------------------
\chapter{Tensores: La Estrella del Data Science}

En inteligencia artificial, los datos rara vez son simples números aislados. Una imagen en color tiene alto, ancho y canales; un lote de textos tiene secuencias de palabras representadas como vectores; un video añade la dimensión temporal. Para manejar estas estructuras de manera eficiente, necesitamos los \textbf{tensores}: la generalización de matrices a más de dos dimensiones. En esta sesión, dominarás su definición, notación, operaciones fundamentales y, lo más importante, cómo se usan en aplicaciones reales de IA.

\section{¿Qué es un tensor?}

Un \textbf{tensor} es un arreglo multidimensional de números. Su \textbf{orden} (o rango) es el número de dimensiones. En IA, trabajamos constantemente con tensores de distintos órdenes:

\begin{itemize}
    \item \textbf{Tensor de orden 0}: escalar (un número). Notación: $x \in \mathbb{R}$.
    \item \textbf{Tensor de orden 1}: vector. Notación: $\mathbf{x} \in \mathbb{R}^{n}$.
    \item \textbf{Tensor de orden 2}: matriz. Notación: $\mathbf{X} \in \mathbb{R}^{m \times n}$.
    \item \textbf{Tensor de orden 3}: por ejemplo, una imagen RGB (alto, ancho, canales). Notación: $\mathcal{X} \in \mathbb{R}^{h \times w \times c}$.
    \item \textbf{Tensor de orden 4}: un lote de imágenes (batch, alto, ancho, canales). Notación: $\mathcal{X} \in \mathbb{R}^{b \times h \times w \times c}$.
    \item \textbf{Tensor de orden 5}: un lote de videos (batch, frames, alto, ancho, canales). Notación: $\mathcal{X} \in \mathbb{R}^{b \times f \times h \times w \times c}$.
\end{itemize}

En general, un tensor de orden $k$ se denota como $\mathcal{X} \in \mathbb{R}^{d_1 \times d_2 \times \dots \times d_k}$, donde $d_i$ es el tamaño de la i-ésima dimensión.

\subsection{Representación visual}

\begin{center}
\begin{tikzpicture}[scale=1.2]
% Escalar
\draw[fill=blue!20] (0,0) circle (0.4);
\node at (0,0) {$x$};
\node at (0,-0.8) {Escalar (0D)};

% Vector
\draw[fill=green!20] (2,-0.5) rectangle (2.5,0.5);
\foreach \y in {-0.4,-0.2,0,0.2,0.4} {
    \fill (2.25,\y) circle (0.05);
}
\node at (2.25,0.8) {Vector (1D)};

% Matriz
\begin{scope}[xshift=4cm, yshift=-0.2cm]
\draw[fill=red!20] (0,0) grid (2,2);
\foreach \i in {0.5,1.5} {
    \foreach \j in {0.5,1.5} {
        \node at (\i,\j) {$x$};
    }
}
\node at (1,2.3) {Matriz (2D)};
\end{scope}

% Tensor 3D
\begin{scope}[xshift=7cm, yshift=0cm]
\draw[fill=yellow!20] (0,0) -- (1,0.5) -- (1,1.5) -- (0,1) -- cycle;
\draw[fill=yellow!40] (0,0) -- (0,1) -- (1,1.5) -- (1,0.5) -- cycle;
\draw[fill=yellow!60] (0,1) -- (1,1.5) -- (2,1) -- (1,0.5) -- cycle;
\node at (0.5,1.8) {Tensor 3D};
\end{scope}
\end{tikzpicture}
\captionof{figure}{Representación de tensores de orden 0 a 3.}
\end{center}

\section{Notación matemática en papers de IA}

En artículos científicos, encontrarás notación como:
\begin{itemize}
    \item $\mathbf{X} \in \mathbb{R}^{n \times d}$: una matriz de $n$ muestras con $d$ características.
    \item $\mathcal{T} \in \mathbb{R}^{b \times c \times h \times w}$: un tensor de lotes (batch size $b$), canales $c$, alto $h$, ancho $w$.
    \item $x_i$: el i-ésimo elemento de un vector.
    \item $X_{ij}$: elemento en la fila $i$, columna $j$ de una matriz.
    \item $\mathcal{T}_{ijk}$ o $\mathcal{T}_{i,j,k}$: elemento de un tensor de orden 3.
\end{itemize}

Es crucial acostumbrarse a esta notación porque es el lenguaje universal de la IA.

\section{Operaciones fundamentales con tensores}

\subsection{Transposición}
La transposición intercambia dimensiones. Para una matriz:
\begin{equation*}
(\mathbf{A}^\top)_{ij} = A_{ji}
\end{equation*}
Para un tensor de orden superior, podemos permutar dimensiones. Por ejemplo, si tenemos un tensor de imágenes $\mathcal{X} \in \mathbb{R}^{b \times h \times w \times c}$ y queremos cambiar el orden de los canales, podemos permutar a $\mathbb{R}^{b \times c \times h \times w}$ (común en PyTorch).

\subsection{Redimensionamiento (Reshape)}
Cambiar la forma de un tensor sin alterar los datos. Por ejemplo, una imagen de $28\times28$ puede aplanarse a un vector de $784$ elementos.
\begin{equation*}
\text{reshape}: \mathbb{R}^{h \times w} \rightarrow \mathbb{R}^{hw}
\end{equation*}

\subsection{Normas de vectores}
Una \textbf{norma} es una función que asigna a un vector su longitud o magnitud. Las más usadas son:

\begin{itemize}
    \item \textbf{Norma L2 o euclidiana}:
    \begin{equation*}
    \|\mathbf{x}\|_2 = \sqrt{\sum_{i=1}^{n} x_i^2}
    \end{equation*}
    Mide la distancia euclidiana desde el origen. Es la norma natural.

    \item \textbf{Norma L1 o de Manhattan}:
    \begin{equation*}
    \|\mathbf{x}\|_1 = \sum_{i=1}^{n} |x_i|
    \end{equation*}
    Suma de valores absolutos. Es más robusta a outliers.

    \item \textbf{Norma L-infinito}:
    \begin{equation*}
    \|\mathbf{x}\|_\infty = \max_i |x_i|
    \end{equation*}
    El máximo valor absoluto.
\end{itemize}

\begin{center}
\begin{tikzpicture}
\begin{axis}[
    title={Círculos unitarios para distintas normas},
    xlabel=$x_1$,
    ylabel=$x_2$,
    axis equal,
    xmin=-1.5, xmax=1.5,
    ymin=-1.5, ymax=1.5,
    grid=major,
    legend pos=north east
]
\addplot[blue, thick, domain=0:90, samples=50] ({cos(x)}, {sin(x)});
\addlegendentry{Norma L2};
\addplot[red, thick] coordinates {(-1,0) (0,1) (1,0) (0,-1) (-1,0)};
\addlegendentry{Norma L1};
\addplot[green, thick] coordinates {(-1,-1) (1,-1) (1,1) (-1,1) (-1,-1)};
\addlegendentry{Norma L-inf};
\end{axis}
\end{tikzpicture}
\captionof{figure}{Conjuntos $\{\mathbf{x} : \|\mathbf{x}\|_p = 1\}$ para p=1,2,$\infty$.}
\end{center}

\section{Aplicaciones en IA y conexión con cursos futuros}

\subsection{Visión por Computador (CV)}
\begin{itemize}
    \item Una \textbf{imagen en escala de grises} es una matriz $\in \mathbb{R}^{h \times w}$ (cada elemento es un píxel entre 0 y 255).
    \item Una \textbf{imagen a color} (RGB) es un tensor $\in \mathbb{R}^{h \times w \times 3}$ (tres canales).
    \item Un \textbf{lote de imágenes} (para entrenar por batches) es un tensor $\in \mathbb{R}^{b \times h \times w \times c}$.
    \item Un \textbf{video} es una secuencia de frames: $\in \mathbb{R}^{f \times h \times w \times c}$, y con batch: $\in \mathbb{R}^{b \times f \times h \times w \times c}$.
\end{itemize}

\subsection{Procesamiento de Lenguaje Natural (NLP)}
\begin{itemize}
    \item Un \textbf{embedding de palabra} es un vector $\in \mathbb{R}^{d}$ (por ejemplo, GloVe, word2vec).
    \item Una \textbf{frase} (secuencia de palabras) es una matriz $\in \mathbb{R}^{n \times d}$, donde $n$ es la longitud de la secuencia.
    \item Un \textbf{lote de frases} (con padding) es un tensor $\in \mathbb{R}^{b \times n \times d}$.
    \item En transformers, la atención usa tensores de dimensiones (batch, heads, secuencia, secuencia).
\end{itemize}

\subsection{Machine Learning (Regularización)}
En regresión lineal y modelos avanzados, se añade un término de regularización a la función de pérdida para evitar sobreajuste. La norma del vector de pesos $\mathbf{w}$ se usa:

\begin{itemize}
    \item \textbf{Regresión Ridge} (regularización L2):
    \begin{equation*}
    \mathcal{L} = \text{MSE} + \lambda \|\mathbf{w}\|_2^2
    \end{equation*}
    Favorece pesos pequeños pero no nulos.

    \item \textbf{Regresión Lasso} (regularización L1):
    \begin{equation*}
    \mathcal{L} = \text{MSE} + \lambda \|\mathbf{w}\|_1
    \end{equation*}
    Favorece pesos dispersos (muchos ceros), útil para selección de características.
\end{itemize}

\subsection{Otros usos}
\begin{itemize}
    \item En \textbf{aprendizaje por refuerzo}, los estados y acciones se representan como vectores.
    \item En \textbf{grafos}, las matrices de adyacencia son tensores.
    \item En \textbf{series temporales}, los datos son tensores 2D (tiempo, características) o 3D (batch, tiempo, características).
\end{itemize}

\section{Python con NumPy: Manipulación de tensores}

NumPy es la librería fundamental para trabajar con tensores en Python. A continuación, ejemplos paso a paso.

\subsection{Creación de tensores}

\begin{lstlisting}[language=Python]
import numpy as np

# Escalar (0D)
escalar = np.array(5)
print(f"Escalar: {escalar}, dimensiones: {escalar.ndim}, forma: {escalar.shape}")

# Vector (1D)
vector = np.array([1, 2, 3, 4])
print(f"Vector: {vector}, ndim: {vector.ndim}, shape: {vector.shape}")

# Matriz (2D)
matriz = np.array([[1, 2], [3, 4], [5, 6]])
print(f"Matriz:\n{matriz}, ndim: {matriz.ndim}, shape: {matriz.shape}")

# Tensor 3D (por ejemplo, 2 imágenes de 3x3 en escala de grises)
tensor_3d = np.random.rand(2, 3, 3)  # 2 imágenes, 3x3
print(f"Tensor 3D shape: {tensor_3d.shape}")

# Tensor 4D (lote de imágenes RGB: 2 imágenes de 32x32 con 3 canales)
tensor_4d = np.random.rand(2, 32, 32, 3)
print(f"Tensor 4D shape: {tensor_4d.shape}")
\end{lstlisting}

\subsection{Atributos importantes}
\begin{lstlisting}[language=Python]
a = np.array([[1,2,3],[4,5,6]])
print(f"Forma: {a.shape}")      # (2,3)
print(f"Nº dimensiones: {a.ndim}")   # 2
print(f"Tamaño total: {a.size}")     # 6
print(f"Tipo de datos: {a.dtype}")   # int64
\end{lstlisting}

\subsection{Transposición}
\begin{lstlisting}[language=Python]
matriz = np.array([[1,2],[3,4],[5,6]])
print("Original:\n", matriz)
print("Transpuesta:\n", matriz.T)

# Para tensores de más dimensiones, usamos transpose con ejes
tensor = np.random.rand(2,3,4)
tensor_permutado = tensor.transpose(2,0,1)  # nuevo orden: (4,2,3)
print(f"Shape original: {tensor.shape}, nuevo: {tensor_permutado.shape}")
\end{lstlisting}

\subsection{Redimensionamiento (reshape)}
\begin{lstlisting}[language=Python]
vector = np.array([1,2,3,4,5,6])
matriz = vector.reshape(2,3)
print("Vector a matriz (2x3):\n", matriz)

# Cuidado: reshape debe ser compatible con el número total de elementos
try:
    mala = vector.reshape(2,4)  # error: 6 elementos no caben en 2x4=8
except ValueError as e:
    print("Error:", e)

# -1 para que NumPy infiera la dimensión
otra = vector.reshape(3, -1)  # 3 filas, columnas = 2
print(otra)
\end{lstlisting}

\subsection{Cálculo de normas}
\begin{lstlisting}[language=Python]
v = np.array([3, 4])

# Norma L2
norma_l2 = np.linalg.norm(v, ord=2)  # por defecto es L2
print(f"Norma L2 de {v}: {norma_l2:.2f}")  # 5.0

# Norma L1
norma_l1 = np.linalg.norm(v, ord=1)
print(f"Norma L1: {norma_l1:.2f}")  # 7.0

# Norma L-inf
norma_inf = np.linalg.norm(v, ord=np.inf)
print(f"Norma infinita: {norma_inf:.2f}")  # 4.0

# Norma L2 al cuadrado (útil para Ridge)
norma_l2_cuad = np.sum(v**2)
print(f"L2^2: {norma_l2_cuad}")  # 25

# Norma por filas en una matriz
matriz = np.array([[1,2],[3,4],[5,6]])
normas_filas = np.linalg.norm(matriz, axis=1)  # norma L2 de cada fila
print("Normas por fila:", normas_filas)
\end{lstlisting}

\section{Ejemplo integrado: Procesamiento de un lote de imágenes}

Simulemos un lote de 3 imágenes en escala de grises de 4x4 píxeles. Calcularemos la norma L2 de cada imagen (como medida de energía) y aplicaremos una normalización simple.

\begin{lstlisting}[language=Python]
import numpy as np

# Crear un lote de 3 imágenes 4x4 (valores aleatorios entre 0 y 255)
batch = np.random.randint(0, 256, size=(3, 4, 4)).astype(np.float32)
print(f"Forma del lote: {batch.shape}")

# Calcular la norma L2 de cada imagen (aplanamos cada imagen a vector)
normas = np.linalg.norm(batch.reshape(3, -1), axis=1)
print("Normas L2 de cada imagen:", normas)

# Normalizar cada imagen dividiendo por su norma (para que tengan norma 1)
batch_normalizado = batch / normas[:, np.newaxis, np.newaxis]
print("Nuevas normas (deberían ser 1):", 
      np.linalg.norm(batch_normalizado.reshape(3,-1), axis=1))

# Si quisiéramos usar regularización L2 en un modelo, el término sería la suma de cuadrados de los pesos.
# Por ejemplo, para una capa con pesos W de forma (features, neuronas), la norma L2 sería:
W = np.random.randn(10, 5)  # 10 características, 5 neuronas
l2_reg = np.sum(W**2)
print(f"Término de regularización L2: {l2_reg:.4f}")
\end{lstlisting}

\section{Conexión con papers científicos: Interpretando notación}

En un artículo de IA, te encontrarás con expresiones como:

\begin{equation*}
\mathcal{L}(\mathbf{w}) = \frac{1}{N} \sum_{i=1}^{N} \| \mathbf{y}_i - \hat{\mathbf{y}}_i \|_2^2 + \lambda \| \mathbf{w} \|_1
\end{equation*}

Interpretación:
\begin{itemize}
    \item $\mathbf{y}_i$ es el vector de etiquetas reales para la i-ésima muestra (podría ser un vector en clasificación multiclase).
    \item $\hat{\mathbf{y}}_i$ es la predicción del modelo.
    \item $\| \cdot \|_2^2$ es la norma L2 al cuadrado (suma de cuadrados).
    \item $\lambda$ es el hiperparámetro de regularización.
    \item $\| \mathbf{w} \|_1$ es la norma L1 del vector de pesos (regularización Lasso).
\end{itemize}

Otro ejemplo típico en transformers:

\begin{equation*}
\text{Attention}(Q,K,V) = \text{softmax}\left( \frac{Q K^\top}{\sqrt{d_k}} \right) V
\end{equation*}

Aquí $Q$, $K$, $V$ son tensores de dimensión (batch, seq_len, d_model). El producto $QK^\top$ implica multiplicación de matrices y produce un tensor de atención.

\section{Conclusión}

Los tensores son el lenguaje de los datos en IA. Dominar su manipulación y las operaciones asociadas (normas, transposición, reshaping) es fundamental para entender cómo se procesan las imágenes, textos, audio y cualquier tipo de dato. En las próximas sesiones, profundizaremos en operaciones más avanzadas y su implementación eficiente.

\end{document}